% Nouveau EGA V (cote 25): archive cote 25, site page 22 (PDF page 23).
% typescript. Draft compared: none.
% AI-transcribed by Claude Opus 5.5 (claude-opus-5-5, adaptive thinking, effort high), Message Batches API, batch ega25b (msgbatch_01A7n5YuXubUPNK1zsepi8RG), 2026-11-13; not yet checked by eye.
\scanpage{25/22}
\mspage{\hand{19}}

5. \uline{Th\'eor\`emes du type Seidenberg}\xout{\unsure{z}}

5.1.
Dans le pr\'esent n\textsuperscript{o}, nous donnons des conditions .... (cf page pr\'ec\'eden-%
te, le texte barr\'e par erreur).

5.2. Comme nous avons vu au n\textsuperscript{o}1, \underline{Y} est d\'efini dans $X x_S \hand{\check{P}} = X_{\check{\underline{P}}}$
comme le ``diviseur'' d'une section $\phi$ de $\underline{O}_X(1)\otimes_S \underline{O}_{\check{\underline{P}}}(1)$, qui induit pour
tout $\hand{\zeta} \in \hand{\check{P}}$ une sedtion $\phi_{\hand{\zeta}}$ de \xout{$\otimes$xxxxx(1)$\otimes$xxxx xxxxxxxxxxxxxx k( )} $\underline{O}_X(1)\otimes_{k(s)}\underline{O}_{\check{\underline{P}}}(1)(\hand{\zeta})$
(faisceau d'ailleurs isomorphe non canoniquement \`a $\underline{O}_X(1)\otimes_{k(s)}k(\hand{\zeta}) =$
$\underline{O}_{X_{k(\hand{\zeta})}}(1)$ ), de sorte que $Y_{\hand{\zeta}}$ n'est autre que le \xout{xxxxx} ``diviseur'' de cet
te section (NB le diviseur d'une section $\phi$ d'un module inversible \underline{L} \margin{$/\!/$}
est d\'efini comme le sous-pr\'esch\'ema ferm\'e d\'efini par l'Id\'eal image de
$\phi^{-1} : \underline{L}^{-1} \longrightarrow \underline{O}$ ). Si \underline{F} est un Module sur X, alors son image inverse \underline{G}
sur \underline{Y} , i.e. l'image inverse de $\underline{F}\otimes_{\underline{O}_S}\underline{O}_{\check{\underline{P}}} = \underline{F}_{\check{\underline{P}}}$ sur le sous-pr\'esch\'ema \underline{Y} de
$X_{\check{\underline{P}}}$ , n'est autre que le conoyau du homomorphisme
\[ \phi^{-1}\otimes \mathrm{id}_{\underline{F}_{\check{\underline{P}}}} : \underline{F}_{\check{\underline{P}}}(-1,-1) \longrightarrow F_{\check{\underline{P}}} , \]
o\`u la notation (-1,-1) s'explique d'elle-m\xout{e}\^eme, comme dirait Mike. De
m\^eme, $G_{\hand{\zeta}}$ est le conayau de l'homomorphisme analogue
\[ F_{k(\hand{\zeta})}(-1\over{)}{,}-1) \longrightarrow F_{k(\hand{\zeta})} \]
o\`u $\hand{\zeta}$ est un point de $\hand{\check{P}}$ (et bien entendu on a une interpr\'etation corres-%
pondante lorsque $\hand{\zeta}$, au lieu d'\^etre un point de $\hand{\check{P}}$ , d\'esigne un\xout{e} point
de $\hand{\check{P}}$ \`a valeurs dans un S' sur S \xout{\unsure{??}}. ). \hand{\unsure{(}}De fa\c{c}on g\'en\'erale, si L est un
Module inversible quelque part, $\phi$ une section, \unsure{d\'afinissant|d\'efinissant} le sous-pr\'esch
ma V($\phi$), alors pour tout Module F , l'image inverse de F sur V($\phi$)
s'identifie, moyennant l'abus de langage habituel, au conayau de $\mathrm{id}_F\otimes\phi$
$F\otimes L^{-1} \longrightarrow F$ ; On dit que $\phi$ est F-\uline{r\'egulier} si l\xout{e}'homomorphisme pr\'ec\'edent
est injectif. Si on choisit un isomorphisme de F et $\underline{O}_X$ , - ce qui
est possible localement - de sorte que $\phi$ s'identifie \`a une section de
$\underline{O}_X$ , cette terminologie est\unsure{b} compatible avec celle \inserted{d\'ej\`a} introduite par ailleur

\margin{22}
